\begin{table}[!ht]
\centering
\caption{Episode marginals for the inflation--expectations system (one-in/one-out)}
\label{tab:h6marginal}
\setlength{\tabcolsep}{6pt}\footnotesize
\begin{threeparttable}
\begin{tabular}{l c c c c}
\toprule
Episode & $\Delta$Overall & $\Delta$Contemp. & $\Delta$Lagged & Lagged \% \\
\midrule
\multicolumn{5}{l}{\textit{$\tau=0.7$}}\\
Great Inflation (pre-1983 in) & +28.5 & +4.8 & +23.7 & 83.1 \\
COVID (2020-- in) & -0.1 & -4.5 & +4.4 & -4533.3 \\
\addlinespace
\multicolumn{5}{l}{\textit{$\tau=0.9$}}\\
Great Inflation (pre-1983 in) & +14.4 & -0.8 & +15.1 & 105.2 \\
COVID (2020-- in) & +12.2 & +3.6 & +8.7 & 70.8 \\
\addlinespace
\bottomrule
\end{tabular}
\begin{tablenotes}[flushleft]\footnotesize
\item \textit{Notes:} One-in/one-out episode marginals for the quarterly two-variable system of Michigan inflation expectations (\texttt{MICH\_QTR}) and quarter-on-quarter CPI inflation, constructed exactly as in Table~\ref{tab:h3}: each row is total connectedness estimated with that episode included (and the other excluded) minus the calm 1983--2019 core, split into contemporaneous and lagged; $n_{\text{lag}}=2$ (a six-month lag horizon). The pre-1983 sample runs from 1960, so the ``Great Inflation'' marginal includes the 1960s run-up. Dependence, not identified causation.
\end{tablenotes}
\end{threeparttable}
\end{table}

